\chapter{INTRODUÇÃO}

\section{Contextualização}

A geração procedural de conteúdo, conhecida pela sigla PCG 
(\textit{Procedural Content Generation}), consiste na utilização de métodos 
computacionais para produzir conteúdos destinados a jogos digitais. 
Sua aplicação abrange diferentes elementos, como mapas, fases, cenários, 
itens e regras, oferecendo possibilidades de criação e variação de conteúdos 
sem que cada elemento precise ser elaborado individualmente. 
Hendrikx et al. (2013) apresentam uma classificação dos tipos de conteúdo 
que podem ser gerados por essas técnicas, enquanto Maleki e Zhao (2024) 
discutem diferentes abordagens empregadas em jogos, incluindo métodos baseados 
em busca, aprendizado de máquina e procedimentos construtivos 
\cite{Hendrikx2013,MalekiZhao2024}.

Entre as aplicações da PCG, a geração de fases para jogos de plataforma 2D 
apresenta particularidades relacionadas à organização do cenário e às 
possibilidades de movimentação do personagem. A disposição das plataformas, 
os desníveis, os vãos e os obstáculos influencia os caminhos disponíveis e os 
desafios enfrentados durante a partida. Nesse contexto, produzir uma fase 
estruturalmente válida não implica, necessariamente, produzir uma fase que 
possa ser completada pelo jogador. Por esse motivo, além da construção dos 
cenários, torna-se importante verificar sua funcionalidade segundo as regras 
e a física estabelecidas para o jogo \cite{SabrawChurchill2022,SchaaBarriga2024}.

Diferentes estratégias podem ser empregadas na construção automática dessas 
fases. A geração aleatória seleciona características do cenário por mecanismos 
de sorteio; a geração baseada em regras utiliza condições explícitas para 
orientar as decisões de construção; e os algoritmos genéticos exploram 
soluções candidatas ao longo de um processo de busca evolutiva. A literatura 
sobre geração procedural baseada em busca destaca a importância da 
representação do conteúdo e das funções de avaliação para orientar a produção 
de soluções \cite{Togelius2011,TogeliusSearch2016}. Sabraw e Churchill (2022), 
por exemplo, utilizam algoritmos genéticos para gerar fases de plataforma do 
tipo \textit{auto-play}, demonstrando uma aplicação da evolução computacional 
nesse domínio \cite{SabrawChurchill2022}.

Entretanto, a comparação entre estratégias de geração não deve se limitar à 
quantidade de fases produzidas ou à obtenção de uma medida isolada de 
qualidade. Schaa e Barriga (2024) analisam o espaço expressivo de diferentes 
geradores de fases de \textit{Super Mario Bros.}, mostrando a utilidade de 
métricas estruturais para caracterizar a variedade dos resultados. Withington, 
Cook e Tokarchuk (2024) identificam diferentes práticas de avaliação de 
geradores de níveis e discutem as dificuldades de estabelecer comparações 
consistentes entre sistemas. De forma complementar, Khalifa et al. (2025) 
propõem um ambiente de referência para avaliar algoritmos de PCG a partir 
de critérios como qualidade, diversidade e controlabilidade 
\cite{SchaaBarriga2024,Withington2024,Khalifa2025}.

A análise computacional também não representa integralmente a experiência do 
jogador. Summerville et al. (2017) investigam relações entre métricas de fases 
de plataforma e avaliações humanas, evidenciando a pertinência de considerar 
tanto características mensuráveis do cenário quanto percepções dos
participantes \cite{Summerville2017}. Assim, a avaliação de estratégias de 
geração pode combinar medidas de completabilidade, diversidade estrutural e 
desempenho computacional com informações obtidas durante a interação dos 
jogadores com as fases produzidas.

Diante desse cenário, este projeto propõe implementar e comparar três 
estratégias de geração procedural em um mesmo protótipo de jogo de plataforma 
2D. A utilização de uma representação comum das fases, de um mesmo modelo 
físico e de procedimentos padronizados de avaliação busca permitir a análise 
das diferenças observadas entre as estratégias, sem pressupor previamente a 
superioridade de qualquer uma delas.

\section{Problema de Pesquisa}

Embora existam diferentes métodos para gerar fases automaticamente, suas 
características e seus resultados dependem das decisões de representação, 
construção e avaliação adotadas em cada sistema. Além disso, fases que atendem 
a restrições estruturais podem não ser completáveis, e medidas computacionais 
não descrevem, isoladamente, a experiência dos jogadores. A definição de 
critérios comuns de comparação é, portanto, relevante para compreender os 
resultados obtidos por diferentes estratégias no mesmo domínio 
experimental \cite{Withington2024,Khalifa2025}.

Diante disso, esta pesquisa busca responder à seguinte questão:
\textbf{Quais diferenças são observadas entre estratégias de geração 
aleatória, geração baseada em regras e geração por algoritmos genéticos, 
aplicadas à criação de fases de jogos de plataforma 2D, quanto à 
completabilidade, à diversidade estrutural, ao desempenho computacional e 
à experiência dos jogadores?}

% ----------------------------------------------------------
% 1.3 OBJETIVOS
% ----------------------------------------------------------

\section{Objetivos}

\subsection{Objetivo Geral}

Analisar comparativamente estratégias de geração aleatória, geração baseada 
em regras e geração por algoritmos genéticos na produção de fases para jogos 
de plataforma 2D, considerando a completabilidade, a diversidade estrutural, 
o desempenho computacional e a experiência dos jogadores.

\subsection{Objetivos Específicos}

\begin{enumerate}
    \item Desenvolver um protótipo de jogo de plataforma 2D com representação 
    computacional padronizada das fases e modelo compartilhado de movimentação 
    e colisão;

    \item Implementar as estratégias de geração aleatória, geração baseada em 
    regras e geração por algoritmos genéticos, utilizando a mesma representação 
    das fases;

    \item Implementar e validar um procedimento automatizado de verificação de 
    completabilidade, considerando as regras de movimentação do personagem e as 
    limitações da busca;

    \item Comparar as estratégias por meio de experimentos computacionais, 
    utilizando métricas de completabilidade, diversidade estrutural e 
    desempenho computacional;

    \item Investigar, de forma exploratória, a experiência dos jogadores com 
    fases produzidas pelas três estratégias, considerando registros das sessões 
    e respostas a questionários.
\end{enumerate}

\section{Justificativa}

A relevância desta pesquisa está na necessidade de compreender não apenas como 
fases podem ser geradas automaticamente, mas também como diferentes estratégias 
se comportam quando avaliadas sob condições comparáveis. A literatura demonstra 
que os geradores podem apresentar diferenças quanto às características dos 
conteúdos produzidos e que a avaliação de sistemas de PCG envolve múltiplos 
critérios, cuja escolha influencia a interpretação dos resultados 
\cite{SchaaBarriga2024,Withington2024,Khalifa2025}. Nesse sentido, a comparação 
proposta pretende produzir evidências sobre o comportamento de três estratégias 
em um domínio específico, com procedimentos e limitações explicitados.

Do ponto de vista técnico, a pesquisa possibilitará desenvolver um protótipo 
experimental no qual os geradores compartilhem a mesma representação de fases, 
a mesma física e os mesmos mecanismos de avaliação. Essa organização permite 
examinar tanto as características das fases quanto o custo computacional 
necessário para produzi-las. Além disso, a verificação automatizada de 
completabilidade será tratada separadamente da validade estrutural, evitando 
considerar uma fase jogável apenas porque seus elementos respeitam limites 
geométricos. A investigação da experiência dos participantes complementará as 
métricas computacionais, embora seus resultados sejam interpretados como 
exploratórios e restritos às fases e aos participantes avaliados 
\cite{Summerville2017,Withington2024}.

No âmbito de Sistemas de Informação, o estudo integra desenvolvimento de 
software, implementação de algoritmos, simulação, registro de dados e análise 
comparativa de resultados. Espera-se que seus procedimentos e resultados 
contribuam para a compreensão das possibilidades e limitações das estratégias 
investigadas, além de oferecerem uma base documentada para estudos posteriores 
sobre geração e avaliação de fases de jogos de plataforma 2D.

\section{Estrutura do Trabalho}

Este trabalho está organizado da seguinte forma:
O Capítulo~\ref{cap:referencial} reúne o referencial teórico, abordando 
geração procedural de conteúdo, geração de fases para jogos de plataforma 2D, 
algoritmos genéticos e trabalhos relacionados. O Capítulo~\ref{cap:metodologia} 
descreve a classificação da pesquisa e os procedimentos previstos para 
implementação dos geradores, avaliação computacional e avaliação com 
participantes. Por fim, o Capítulo~\ref{cap:cronograma} apresenta o 
planejamento das atividades para a execução do trabalho.